\documentclass[11pt,a4paper]{article}

\usepackage[T1]{fontenc}
\usepackage[utf8]{inputenc}
\usepackage{lmodern}
\usepackage[margin=1in]{geometry}
\usepackage{amsmath,amssymb,amsthm}
\usepackage{array}

\setlength{\emergencystretch}{1.5em}
\renewcommand{\arraystretch}{1.22}

\newtheorem{theorem}{Theorem}[section]
\newtheorem{lemma}[theorem]{Lemma}
\newtheorem{proposition}[theorem]{Proposition}
\newtheorem{corollary}[theorem]{Corollary}
\theoremstyle{definition}
\newtheorem{definition}[theorem]{Definition}
\newtheorem{remark}[theorem]{Remark}
\newtheorem{example}[theorem]{Example}

\newcommand{\F}{\mathbb{F}}
\newcommand{\Z}{\mathbb{Z}}
\newcommand{\N}{\mathbb{N}}
\newcommand{\SG}{\mathcal{G}}
\newcommand{\mex}{\mathrm{mex}}
\newcommand{\supp}{\mathrm{supp}}
\newcommand{\ev}{\mathrm{ev}}

\title{\textbf{Disproof and Complete Algebraic Resolution of Conjecture 00000001288:\\
Sprague--Grundy Functions of Turning Turtles and\\
Homomorphisms of $\Z_2[x]$}}
\author{Formalized in Core Lean 4.19.0 with Independent Computational Verification}
\date{}

\begin{document}
\maketitle

\begin{abstract}
We resolve Conjecture \texttt{00000001288}, which asserts that the Sprague--Grundy (SG)
function of the impartial coin-turning game \emph{Turning Turtles} is generated by a
homomorphism of the polynomial ring $\Z_2[x]$. Under the standard mathematical meaning
of a homomorphism of the polynomial ring $\Z_2[x]$---namely, a ring (or even non-unital
rng) homomorphism preserving both polynomial addition and multiplication---the conjecture
is \textbf{false}.

After establishing the exact $\mex$-based Sprague--Grundy formulas under both $0$-based
indexing ($\SG_0(\sum_{k\in S} x^k) = \bigoplus_{k\in S}(k+1)$) and $1$-based indexing
($\SG_1(\sum_{k\in S} x^k) = \bigoplus_{k\in S} k$), we uncover both the algebraic origin
of the conjecture and its complete impossibility:
\begin{enumerate}
  \setlength{\itemsep}{2pt}
  \item \textbf{Additive $\F_2$-linearity vs.\ ring multiplicativity:} Both $\SG_0$ and
    $\SG_1$ are surjective $\F_2$-linear (additive group) homomorphisms
    $(\Z_2[x],+) \to (\N,\oplus)$, a property shared by all non-interfering coin-turning
    games that uses only the infinite-dimensional $\F_2$-vector space structure of $\Z_2[x]$.
  \item \textbf{Small-board nimber coincidence ($\deg \le 2$):} Equipping $\N$ with
    Conway's nimber field operations $(\oplus,\otimes)$, the unique unital ring
    homomorphism $\ev_2 : \Z_2[x] \to (\N,\oplus,\otimes)$ sending $x \mapsto 2$ agrees
    identically with $\SG_0$ on all $2^3 = 8$ polynomials of degree $\le 2$
    ($1 \mapsto 1$, $x \mapsto 2$, $x^2 \mapsto 3$), before failing at the fourth coin:
    $\SG_0(x^3) = 4 \neq 1 = 2^{\otimes 3} = \ev_2(x^3)$.
  \item \textbf{Universal ring and quotient obstructions:} Not only does no ring
    homomorphism $\Z_2[x] \to (\N,\oplus,\otimes)$ generate $\SG_0$ or $\SG_1$ (every
    such homomorphism has bounded image inside a finite Fermat subfield
    $[0,2^{2^m}) \cong \F_{2^{2^m}}$), but for \emph{every} ring or non-unital rng
    $(R,\boxplus,\boxtimes)$ and \emph{every} zero-reflecting map $\iota : \N \to R$,
    neither $\iota \circ \SG_0$ on $\Z_2[x]$ nor $\iota \circ \SG_1$ on $x\Z_2[x]$ is a
    rng homomorphism because $\ker(\SG)$ is not an ideal. Stronger still, $\ker(\SG)$
    contains no nonzero ideal of $\Z_2[x]$ (zero conductor), forcing any rng homomorphism
    $\Phi : \Z_2[x] \to R$ through which $\SG = \psi \circ \Phi$ factors via an arbitrary
    set map $\psi : R \to \N$ to be injective ($\ker(\Phi)=\{0\}$).
\end{enumerate}
All definitions and theorems are machine-checked from first principles in Core Lean~4.19.0
with zero external dependencies and zero \texttt{sorry} axioms.
\end{abstract}

\section{Introduction and Problem Statement}

Conjecture \texttt{00000001288} from \emph{The Last Math Competition} reads as follows:

\begin{quote}
\textbf{Conjecture 00000001288.} \emph{The SG function of the Turning Turtles game is
generated by a homomorphism of the polynomial ring over $\Z_2$ (the algebraization of
Turning Turtles).}
\end{quote}

In the classical theory of impartial combinatorial games developed by Berlekamp, Conway,
and Guy in \emph{Winning Ways for Your Mathematical Plays}~\cite{WW}, Chapter~14
(``Turn and Turn About''), \emph{Turning Turtles} is the prototype of one-dimensional
coin-turning games. A configuration consists of a semi-infinite row of coins (or turtles),
finitely many of which show Heads (upright, state $1 \in \Z_2$) while all remaining coins
show Tails (upside-down, state $0 \in \Z_2$). Under normal play, a legal move turns one
Heads coin over to Tails and \emph{optionally} flips at most one strictly preceding coin
(either from Heads to Tails or from Tails to Heads); a player unable to move loses.

\paragraph{Algebraization via $\Z_2[x]$.}
Identifying a Heads coin at position $k$ with the monomial $x^k \in \Z_2[x]$, every game
position with Heads set $S \subset_{\mathrm{fin}} \N$ corresponds bijectively to the
polynomial
\[
  P_S(x) = \sum_{k \in S} x^k \in \Z_2[x].
\]
Because flipping a coin at index $j$ toggles its state modulo $2$, flipping a single coin
at $n$ adds $x^n$ in $\Z_2[x]$, while flipping coins at $n$ and $i < n$ adds
$x^n + x^i = x^i(x^{n-i} + 1)$ in $\Z_2[x]$. Two indexing conventions appear in the
combinatorial game literature:
\begin{itemize}
  \setlength{\itemsep}{2pt}
  \item \textbf{$0$-based convention ($\SG_0$ on $\Z_2[x]$):} Active coins occupy indices
    $n \in \{0,1,2,\dots\}$, putting the entire polynomial ring $\Z_2[x]$ in bijection
    with game positions. Flipping $n$ optionally flips one $i$ with $0 \le i < n$.
  \item \textbf{$1$-based convention ($\SG_1$ on $x\Z_2[x]$ or $\Z_2[x]$, following
    \emph{Winning Ways}~\cite{WW}):} Active coins occupy indices $n \in \{1,2,3,\dots\}$,
    and an optional second flip acts on $1 \le i < n$. Configurations with no coin at $0$
    form the maximal ideal (non-unital subring) $I_+ = x\Z_2[x] = \{P \in \Z_2[x] : P(0)=0\}$,
    or extend to all of $\Z_2[x]$ by treating index $0$ as inert ($\SG_1(1)=0$).
\end{itemize}

\paragraph{Summary of Verdict and Main Results.}
In abstract algebra, a \emph{homomorphism of the polynomial ring $\Z_2[x]$} means a ring
homomorphism out of $(\Z_2[x],+,\cdot)$. Because $\Z_2[x]$ is the free commutative
$\Z_2$-algebra on the single generator $x$, a ring homomorphism $\Phi : \Z_2[x] \to R$
is completely generated by the single element $\Phi(x) \in R$ (together with $\Phi(1)$ in
the non-unital setting), which explains the phrasing ``\emph{generated by} a homomorphism
of the polynomial ring over $\Z_2$.''

We prove that \textbf{Conjecture \texttt{00000001288} is false} under both indexing
conventions and under every ring-theoretic interpretation:
\begin{enumerate}
  \setlength{\itemsep}{2pt}
  \item \textbf{Exact SG Evaluation (Theorem~\ref{thm:sg-exact}):} Directly from the
    recursive $\mex$ definition on the game graph of $\Z_2[x]$,
    $\SG_0(\sum_{k\in S} x^k) = \bigoplus_{k\in S}(k+1)$ and
    $\SG_1(\sum_{k\in S, k\ge 1} x^k) = \bigoplus_{k\in S, k\ge 1} k$.
  \item \textbf{Positive Structure \& Small-Board Coincidence
    (Propositions~\ref{prop:linear}--\ref{prop:coincidence}):} Both $\SG_0$ and $\SG_1$
    are surjective $\F_2$-linear homomorphisms $(\Z_2[x],+) \to (\N,\oplus)$. Moreover,
    into Conway's nimber field $(\N,\oplus,\otimes)$, the unital ring homomorphism
    $\ev_2(P) = \bigoplus_{k\in\supp(P)} 2^{\otimes k}$ coincides with $\SG_0(P)$ for
    \emph{all} $8$ polynomials of degree $\le 2$, failing first at $x^3$ where
    $\SG_0(x^3) = 4 \neq 1 = \ev_2(x^3)$.
  \item \textbf{Complete Ring and Factorization Impossibility
    (Theorems~\ref{thm:nimber-no}--\ref{thm:zero-conductor}):}
    \begin{itemize}
      \setlength{\itemsep}{1pt}
      \item Every ring homomorphism $\Phi : \Z_2[x] \to (\N,\oplus,\otimes)$ has bounded
        image inside a finite Fermat subfield $[0,2^{2^m})$, whereas $\SG_0,\SG_1$ are
        unbounded surjections onto $\N$.
      \item For \emph{any} (not necessarily commutative or unital) ring/rng
        $(R,\boxplus,\boxtimes)$ and \emph{any} map $\iota : \N \to R$ with
        $\iota(n)=\iota(0)\Rightarrow n=0$, neither $\iota \circ \SG_0$ on $\Z_2[x]$ nor
        $\iota \circ \SG_1$ on $x\Z_2[x]$ is a rng homomorphism: explicitly,
        $\SG_0(1+x+x^2)=0$ whereas $\SG_0(x(1+x+x^2))=5\neq 0$, and
        $\SG_1(x+x^2+x^3)=0$ whereas $\SG_1(x(x+x^2+x^3))=5\neq 0$.
      \item For every nonzero $m \in \Z_2[x]$, there exists $k \ge 1$ with
        $\SG_0(x^k m) \neq 0$ and $\SG_1(x^k m) \neq 0$. Hence any rng homomorphism
        $\Phi : \Z_2[x] \to R$ through which $\SG_0$ or $\SG_1$ factors via \emph{any}
        set function $\psi : R \to \N$ ($\SG = \psi \circ \Phi$) has trivial kernel
        $\ker(\Phi)=\{0\}$ and is injective.
    \end{itemize}
\end{enumerate}

\section{Game Formalization and Exact Sprague--Grundy Formulas}

For any polynomial $P(x) = \sum_{k\ge 0} c_k x^k \in \Z_2[x]$ ($c_k \in \{0,1\}$, finitely
many $c_k = 1$), write $\supp(P) = \{k \in \N : c_k = 1\}$ and define its binary weight
valuation $\nu(P) = \sum_{k\in\supp(P)} 2^k \in \N$. Note that $P \mapsto \nu(P)$ is a
bijection $\Z_2[x] \xrightarrow{\sim} \N$ with $\nu(P + Q) = \nu(P) \oplus \nu(Q)$, where
$\oplus$ denotes bitwise exclusive-or (nim-addition) on $\N$.

\begin{definition}[Turning Turtles Move Relations and SG Functions]
\label{def:tt}
On the position space $\Z_2[x]$:
\begin{enumerate}
  \setlength{\itemsep}{2pt}
  \item \textbf{$0$-indexed Turning Turtles ($\to_0$):} $P \to_0 Q$ holds if and only if
    there exists $n \in \supp(P)$ such that either $Q = P + x^n$, or there exists
    $0 \le i < n$ with $Q = P + x^n + x^i$.
  \item \textbf{$1$-indexed Turning Turtles ($\to_1$):} $P \to_1 Q$ holds if and only if
    there exists $n \in \supp(P)$ with $n \ge 1$ such that either $Q = P + x^n$, or there
    exists $1 \le i < n$ with $Q = P + x^n + x^i$.
\end{enumerate}
For $t \in \{0,1\}$ and any finite set $A \subset \N$, let
$\mex(A) = \min(\N \setminus A)$. The Sprague--Grundy functions $\SG_0, \SG_1 : \Z_2[x] \to \N$
are defined recursively by
\[
  \SG_t(P) = \mex\{\SG_t(Q) : P \to_t Q\} \qquad (t \in \{0,1\}).
\]
\end{definition}

\begin{lemma}[Strict Termination]
\label{lem:wf}
Whenever $P \to_0 Q$ or $P \to_1 Q$, there are only finitely many options $Q$ from $P$,
and $\nu(Q) < \nu(P)$. Hence $\to_0$ and $\to_1$ are finite, well-founded impartial games,
and $\SG_0, \SG_1$ are uniquely well-defined on all of $\Z_2[x]$.
\end{lemma}

\begin{proof}
If $Q = P + x^n$ with $n \in \supp(P)$, then $\nu(Q) = \nu(P) - 2^n < \nu(P)$. If
$Q = P + x^n + x^i$ with $n \in \supp(P)$ and $0 \le i < n$, then removing $n$ subtracts
$2^n$ while toggling $i$ changes $\nu$ by $\pm 2^i$; since $2^i \le 2^{n-1} < 2^n$, we have
$\nu(Q) \le \nu(P) - 2^n + 2^i < \nu(P)$. Since $\supp(P)$ is finite and $i < n$ ranges
over a finite set for each $n \in \supp(P)$, each position has finitely many options.
\end{proof}

Notice that when $i \in \supp(P)$ already holds, the move $P \mapsto P + x^n + x^i$ turns
both coin $n$ and coin $i$ from Heads to Tails ($c_i : 1 \mapsto 0$). Rather than appealing
informally to cancellation of duplicate coins in a disjunctive sum, we establish the exact
formula directly from the $\mex$ characterization on $\Z_2[x]$.

\begin{theorem}[Exact Sprague--Grundy Formulas for Turning Turtles]
\label{thm:sg-exact}
For every polynomial $P(x) = \sum_{k \in S} x^k \in \Z_2[x]$:
\begin{equation}
\label{eq:sg-formulas}
  \SG_0(P) = \bigoplus_{k \in S} (k + 1), \qquad
  \SG_1(P) = \bigoplus_{\substack{k \in S \\ k \ge 1}} k = \bigoplus_{k \in S} k.
\end{equation}
In particular, $\SG_0(x^n) = n + 1$ for all $n \ge 0$, while $\SG_1(1) = 0$ and
$\SG_1(x^n) = n$ for all $n \ge 1$.
\end{theorem}

\begin{proof}
Write $W_0(P) = \bigoplus_{k \in \supp(P)} (k+1)$ and $W_1(P) = \bigoplus_{k \in \supp(P)} k$.
Note first that $\supp(P + Q) = \supp(P) \triangle \supp(Q)$, and since $a \oplus a = 0$
for all $a \in \N$, symmetric-difference cancellation yields
$W_t(P + Q) = W_t(P) \oplus W_t(Q)$ for all $P,Q \in \Z_2[x]$ and $t \in \{0,1\}$.

We prove $\SG_t(P) = W_t(P)$ by strong induction on $\nu(P)$. Assume $\SG_t(Q) = W_t(Q)$
for all $Q$ with $\nu(Q) < \nu(P)$. By Lemma~\ref{lem:wf} and the definition of $\mex$, it
suffices to verify: (i) $W_t(Q) \neq W_t(P)$ whenever $P \to_t Q$, and (ii) for every
integer $m$ with $0 \le m < W_t(P)$, there exists $Q$ with $P \to_t Q$ and $W_t(Q) = m$.

\smallskip
\noindent\emph{Step (i): Excludant property.}
For $t = 0$: if $Q = P + x^n$, then $W_0(Q) = W_0(P) \oplus (n+1) \neq W_0(P)$ since
$n + 1 \ge 1 \neq 0$; if $Q = P + x^n + x^i$ with $0 \le i < n$, then
$W_0(Q) = W_0(P) \oplus (n+1) \oplus (i+1) \neq W_0(P)$ since $i + 1 \neq n + 1$ implies
$(n+1) \oplus (i+1) \neq 0$. Similarly, for $t = 1$: if $Q = P + x^n$ with $n \ge 1$, then
$W_1(Q) = W_1(P) \oplus n \neq W_1(P)$ since $n \ge 1$; if $Q = P + x^n + x^i$ with
$1 \le i < n$, then $W_1(Q) = W_1(P) \oplus n \oplus i \neq W_1(P)$ since $i \neq n$.

\smallskip
\noindent\emph{Step (ii): Minimality property.}
Let $0 \le m < g := W_t(P)$, and set $d = g \oplus m > 0$. Let $b \ge 0$ be the most
significant bit at which the binary expansions of $g$ and $m$ differ; since $m < g$, bit $b$
of $g$ is $1$, bit $b$ of $m$ is $0$, and for all $j > b$ bit $j$ of $d = g \oplus m$ is
$0$ (while bit $b$ of $d$ is $1$).

\emph{Case $t = 0$:} Since bit $b$ of $W_0(P) = \bigoplus_{k\in\supp(P)} (k+1)$ is $1$,
there exists $n \in \supp(P)$ such that bit $b$ of $n + 1$ is $1$. Consider
$v = (n + 1) \oplus d$. For every $j > b$, bit $j$ of $d$ is $0$, so $v$ and $n+1$ agree
above bit $b$; at bit $b$, both $n+1$ and $d$ have bit $1$, so $v$ has bit $0$. Hence
$0 \le v < n + 1$.
\begin{itemize}
  \setlength{\itemsep}{2pt}
  \item If $v = 0$, then $n + 1 = d$. Choosing $Q = P + x^n$ gives $P \to_0 Q$ and
    $W_0(Q) = g \oplus (n+1) = g \oplus d = m$.
  \item If $v \ge 1$, write $v = i + 1$ with $0 \le i < n$. Choosing $Q = P + x^n + x^i$
    gives $P \to_0 Q$ and
    $W_0(Q) = g \oplus (n+1) \oplus (i+1) = g \oplus (n+1) \oplus ((n+1) \oplus d)
    = g \oplus d = m$.
\end{itemize}

\emph{Case $t = 1$:} Since bit $b$ of $W_1(P) = \bigoplus_{k\in\supp(P)} k$ is $1$, there
exists $n \in \supp(P)$ such that bit $b$ of $n$ is $1$; in particular $n \ge 1$. Setting
$v = n \oplus d$, the same highest-differing-bit comparison yields $0 \le v < n$. If
$v = 0$, then $n = d$ and $Q = P + x^n$ satisfies $P \to_1 Q$ and $W_1(Q) = g \oplus n = m$;
if $v \ge 1$, then $1 \le v < n$ and $Q = P + x^n + x^v$ satisfies $P \to_1 Q$ and
$W_1(Q) = g \oplus n \oplus v = g \oplus d = m$.
\end{proof}

\section{Positive Structure and the Small-Board Nimber Coincidence}

Before proving the impossibility theorems, we isolate the two genuine algebraic phenomena
that motivate Conjecture~\texttt{00000001288}.

\subsection{Surjective $\F_2$-Linearity on $(\Z_2[x],+)$}

\begin{proposition}[Additive / $\F_2$-Linear Homomorphism]
\label{prop:linear}
Both $\SG_0 : (\Z_2[x],+,0) \to (\N,\oplus,0)$ and
$\SG_1 : (\Z_2[x],+,0) \to (\N,\oplus,0)$ (as well as the restriction
$\SG_1|_{x\Z_2[x]} : (x\Z_2[x],+,0) \to (\N,\oplus,0)$) are surjective homomorphisms of
$\F_2$-vector spaces (equivalently, of elementary abelian $2$-groups):
\[
  \SG_t(0) = 0, \qquad \SG_t(P + Q) = \SG_t(P) \oplus \SG_t(Q) \quad
  (\forall P,Q \in \Z_2[x],\; t \in \{0,1\}),
\]
with $\SG_0(x^{g-1}) = g$ and $\SG_1(x^g) = g$ for every integer $g \ge 1$.
\end{proposition}

\begin{proof}
Additivity was shown in the proof of Theorem~\ref{thm:sg-exact} via
$\supp(P+Q) = \supp(P)\triangle\supp(Q)$, and scalar compatibility over $\Z_2 = \{0,1\}$
follows immediately from $\SG_t(0) = 0$. For surjectivity, $\SG_0(0) = \SG_1(0) = 0$, and
for $g \ge 1$ we have $\SG_0(x^{g-1}) = g$ while $x^g \in x\Z_2[x]$ satisfies
$\SG_1(x^g) = g$.
\end{proof}

\begin{remark}[Why $\F_2$-linearity is not a polynomial ring homomorphism]
As an $\F_2$-vector space, $\Z_2[x] = \bigoplus_{n\ge 0} \F_2 x^n$ is free on the
\emph{countably infinite} basis $\{1,x,x^2,\dots\}$, so \emph{every} sequence
$(w_n)_{n\ge 0} \in \N^\N$ defines a unique $\F_2$-linear map
$\sum_{k\in S} x^k \mapsto \bigoplus_{k\in S} w_k$. Indeed, every impartial coin-turning
game without cross-coin interference (Mock Turtles, Ruler, Grunt, Twins, etc.~\cite{WW})
induces an $\F_2$-linear SG map by the Sprague--Grundy theorem and modulo-$2$ cancellation
$x^n + x^n = 0$. By contrast, as a \emph{polynomial ring}, $\Z_2[x]$ is generated over
$\Z_2$ by the single indeterminate $x$; a homomorphism of the polynomial ring $\Z_2[x]$
must preserve polynomial multiplication $P \cdot Q$.
\end{remark}

\subsection{Conway's Nimber Field and the Degree-$\le 2$ Coincidence}

Recall that Conway~\cite{ONAG,WW} equips $\N$ with nim-addition $\oplus$ and recursive
nim-multiplication
\begin{equation}
\label{eq:nimmul}
  a \otimes b = \mex\bigl\{(a' \otimes b) \oplus (a \otimes b') \oplus (a' \otimes b')
  \;\big|\; 0 \le a' < a,\; 0 \le b' < b\bigr\},
\end{equation}
making $(\N,\oplus,\otimes)$ an algebraically closed field of characteristic $2$
isomorphic to $\overline{\F_2}$. Specifically, if $F_m = 2^{2^m}$ ($m \ge 0$) denotes the
$m$-th Fermat $2$-power ($F_0 = 2, F_1 = 4, F_2 = 16, \dots$), Conway's simplicity rules
specify $a \otimes F_m = a \cdot F_m$ for $a < F_m$ and
$F_m \otimes F_m = \frac{3}{2} F_m$, so each initial interval
$[0, F_m) = \{0,1,\dots,2^{2^m}-1\}$ is a finite subfield isomorphic to $\F_{2^{2^m}}$
(with $\F_2 = \{0,1\}$ and $\F_4 = \{0,1,2,3\}$).

By the universal property of $\Z_2[x]$, for every $\alpha \in \N$ there is a unique unital
ring homomorphism $\ev_\alpha : (\Z_2[x],+,\cdot,0,1) \to (\N,\oplus,\otimes,0,1)$ with
$\ev_\alpha(x) = \alpha$, given explicitly by
\[
  \ev_\alpha\Bigl(\sum_{k \in S} x^k\Bigr) = \bigoplus_{k \in S} \alpha^{\otimes k}.
\]
Setting $\alpha = \SG_0(x) = 2 = F_0$ produces a striking small-board coincidence:

\begin{proposition}[Degree-$\le 2$ Coincidence and Degree-$3$ Breakdown]
\label{prop:coincidence}
Let $\ev_2 : \Z_2[x] \to (\N,\oplus,\otimes)$ be the unital ring homomorphism sending
$x \mapsto 2$. Then:
\begin{enumerate}
  \setlength{\itemsep}{2pt}
  \item For \emph{every} polynomial $P \in \Z_2[x]$ with $\deg(P) \le 2$ (all $2^3 = 8$
    configurations on the first three coins $\{0,1,2\}$), $\SG_0(P) = \ev_2(P)$.
  \item At $P(x) = x^3$ (the fourth coin), $\SG_0(x^3) = 4$ whereas $\ev_2(x^3) = 1$.
\end{enumerate}
\end{proposition}

\begin{proof}
Computing directly from~\eqref{eq:nimmul} on $\F_4 = \{0,1,2,3\}$ gives $2^{\otimes 0} = 1$,
$2^{\otimes 1} = 2$, and
\[
  2^{\otimes 2} = 2 \otimes 2
  = \mex\{0 \oplus 0 \oplus 0,\; 0 \oplus 2 \oplus 0,\; 2 \oplus 0 \oplus 0,\;
          2 \oplus 2 \oplus 1\}
  = \mex\{0,2,1\} = 3.
\]
Meanwhile, Theorem~\ref{thm:sg-exact} gives $\SG_0(1) = 1$, $\SG_0(x) = 2$, and
$\SG_0(x^2) = 3$. Because both $\SG_0$ and $\ev_2$ are $\F_2$-linear maps
$(\Z_2[x],+) \to (\N,\oplus)$ agreeing on the basis $\{1,x,x^2\}$, they agree on all $8$
polynomials in $\mathrm{span}_{\F_2}\{1,x,x^2\}$. At $n = 3$, since $2 \in \F_4$ satisfies
$2^{\otimes 2} \oplus 2 \oplus 1 = 3 \oplus 2 \oplus 1 = 0$, we obtain
$2^{\otimes 3} = 2 \otimes 3 = 1$, whereas $\SG_0(x^3) = 3 + 1 = 4 \neq 1$.
\end{proof}

\begin{table}[ht]
\centering
\caption{Comparison of $0$-indexed Turning Turtles $\SG_0(x^n)$ against the nimber
evaluation ring homomorphism $\ev_2(x^n) = 2^{\otimes n}$ for monomials $x^n$ ($0 \le n \le 6$).}
\label{tab:coincidence}
\medskip
\begin{tabular}{cccccccc}
\hline
Monomial $x^n$ & $1$ & $x$ & $x^2$ & $x^3$ & $x^4$ & $x^5$ & $x^6$ \\
\hline
$\SG_0(x^n) = n + 1$ & $1$ & $2$ & $3$ & $\mathbf{4}$ & $5$ & $6$ & $7$ \\
$\ev_2(x^n) = 2^{\otimes n}$ & $1$ & $2$ & $3$ & $\mathbf{1}$ & $2$ & $3$ & $1$ \\
Status & Match & Match & Match & \textbf{Fails} & Fails & Fails & Fails \\
\hline
\end{tabular}
\end{table}

\begin{remark}[Combinatorial interpretation of $\ev_2$]
The periodic sequence $\ev_2(x^n) = 2^{\otimes n} = (1,2,3,1,2,3,\dots)$ is the exact
Sprague--Grundy function of \emph{Bounded-Jump Turning Turtles} with maximum leftward jump
distance $2$, in which flipping coin $n$ optionally flips one coin $i \in \{n-1,n-2\}\cap\N$.
On boards of length $\le 3$ ($n \in \{0,1,2\}$), the constraints $i \in \{n-1,n-2\}\cap\N$
and $0 \le i < n$ define the exact same legal moves; at $n = 3$, full Turning Turtles
additionally permits flipping $(3,0)$ (jump distance $3$), lifting $\SG_0(x^3)$ from $1$
to the next Fermat $2$-power $4 = F_1$.
\end{remark}

\section{Complete Impossibility Under Every Ring Interpretation}

We now prove that no ring structure on the codomain and no quotient-ring factorization can
rescue Conjecture~\texttt{00000001288}.

\subsection{Impossibility in Conway's Nimber Field $(\N,\oplus,\otimes)$}

\begin{theorem}[Refutation in Conway's Nimber Ring]
\label{thm:nimber-no}
Neither $\SG_0$ nor $\SG_1$ is a ring homomorphism (or even a non-unital rng homomorphism)
from $\Z_2[x]$ into Conway's nimber ring $(\N,\oplus,\otimes)$. Moreover, \emph{every}
rng homomorphism $\Phi : \Z_2[x] \to (\N,\oplus,\otimes)$ has finite, bounded image
contained in some Fermat subfield $[0,2^{2^m})$, whereas $\SG_0(\Z_2[x]) = \SG_1(x\Z_2[x]) = \N$.
\end{theorem}

\begin{proof}
For $\SG_0$, multiplicativity $\SG_0(P \cdot Q) = \SG_0(P) \otimes \SG_0(Q)$ fails at
$P = x$, $Q = x^2$:
\[
  \SG_0(x \cdot x^2) = \SG_0(x^3) = 4 \neq 1 = 2 \otimes 3 = \SG_0(x) \otimes \SG_0(x^2).
\]
For $\SG_1$, multiplicativity fails already on $x\Z_2[x]$ at $P = Q = x$:
\[
  \SG_1(x \cdot x) = \SG_1(x^2) = 2 \neq 1 = 1 \otimes 1 = \SG_1(x) \otimes \SG_1(x).
\]
More generally, let $\Phi : \Z_2[x] \to (\N,\oplus,\otimes)$ be any rng homomorphism. Since
$\Phi(1) = \Phi(1 \cdot 1) = \Phi(1) \otimes \Phi(1)$ in the field $(\N,\oplus,\otimes)$,
either $\Phi(1) = 0$ (forcing $\Phi(P) = \Phi(P \cdot 1) = \Phi(P) \otimes 0 = 0$ for all
$P$) or $\Phi(1) = 1$. In the latter case $\Phi = \ev_\alpha$ where $\alpha = \Phi(x)$.
Choosing $m \ge 0$ with $\alpha < F_m = 2^{2^m}$, closure of the Fermat subfield
$[0, F_m) \cong \F_{2^{2^m}}$ under $\oplus$ and $\otimes$ gives $\alpha^{\otimes k} < F_m$
for all $k \ge 0$ and hence $\Phi(P) = \bigoplus_{k\in\supp(P)} \alpha^{\otimes k} < F_m$
for all $P \in \Z_2[x]$. By Proposition~\ref{prop:linear}, $\SG_0$ and $\SG_1$ are
surjective onto $\N$, so neither can equal any such $\Phi$.
\end{proof}

\subsection{Universal Ring/Rng Obstruction for Arbitrary Codomain Rings}

One might ask whether $\N$ (or a superset $R$ encoding $\N$) could be endowed with some
\emph{other} ring operations $(\boxplus,\boxtimes)$---not equal to Conway's $(\oplus,\otimes)$,
and possibly non-commutative or non-unital---under which $\SG_0$ or $\SG_1$ becomes a
homomorphism. Recall that a \emph{rng} $(R,\boxplus,\boxtimes,0_R)$ is an abelian group
$(R,\boxplus,0_R)$ equipped with an associative multiplication $\boxtimes$ distributing
on both sides over $\boxplus$ (every unital ring is in particular a rng).

\begin{theorem}[Universal Ring/Rng Obstruction]
\label{thm:universal-rng}
Let $(R,\boxplus,\boxtimes,0_R)$ be \emph{any} rng, and let $\iota : \N \to R$ be \emph{any}
map that distinguishes nonzero nimbers from $0$, i.e., $\iota(n) = \iota(0) \Rightarrow n = 0$
(in particular, any injective map $\iota : \N \hookrightarrow R$, or $\iota = \mathrm{id}_\N$
for any hypothetical rng structure on $\N$).
\begin{enumerate}
  \setlength{\itemsep}{2pt}
  \item The composite $\iota \circ \SG_0 : \Z_2[x] \to R$ is \textbf{not} a rng homomorphism.
  \item Even restricted to the positive-degree sub-rng $x\Z_2[x]$ (and hence also on all of
    $\Z_2[x]$), the composite $\iota \circ \SG_1 : x\Z_2[x] \to R$ is \textbf{not} a rng
    homomorphism.
\end{enumerate}
\end{theorem}

\begin{proof}
In any rng $(R,\boxplus,\boxtimes,0_R)$, $r \boxtimes 0_R = r \boxtimes (0_R \boxplus 0_R)
= (r \boxtimes 0_R) \boxplus (r \boxtimes 0_R)$, so cancelling $r \boxtimes 0_R$ in the
abelian group $(R,\boxplus,0_R)$ yields $r \boxtimes 0_R = 0_R$ for all $r \in R$. Likewise,
any additive homomorphism $\Phi$ from $\Z_2[x]$ (or $x\Z_2[x]$) to $R$ satisfies
$\Phi(0) = \Phi(0 + 0) = \Phi(0) \boxplus \Phi(0)$, hence $\Phi(0) = 0_R$.

\smallskip
\noindent(1) Suppose $\Phi = \iota \circ \SG_0 : \Z_2[x] \to R$ were a rng homomorphism.
Since $\SG_0(0) = 0$, we have $\iota(0) = \Phi(0) = 0_R$. Consider
$f_0(x) = 1 + x + x^2 \in \Z_2[x]$. By Theorem~\ref{thm:sg-exact},
\[
  \SG_0(f_0) = 1 \oplus 2 \oplus 3 = 0, \qquad
  \SG_0(x \cdot f_0) = \SG_0(x + x^2 + x^3) = 2 \oplus 3 \oplus 4 = 1 \oplus 4 = 5.
\]
Then $\Phi(f_0) = \iota(0) = 0_R$, so multiplicativity of $\Phi$ forces
\[
  \iota\bigl(\SG_0(x \cdot f_0)\bigr) = \Phi(x \cdot f_0) = \Phi(x) \boxtimes \Phi(f_0)
  = \Phi(x) \boxtimes 0_R = 0_R = \iota(0).
\]
By the hypothesis on $\iota$, $\SG_0(x \cdot f_0) = 0$, contradicting $\SG_0(x \cdot f_0) = 5$.

\smallskip
\noindent(2) Suppose $\Phi = \iota \circ \SG_1 : x\Z_2[x] \to R$ were a rng homomorphism on
$x\Z_2[x]$. Again $\iota(0) = \Phi(0) = 0_R$. Consider
$f_1(x) = x + x^2 + x^3 \in x\Z_2[x]$ and $x \in x\Z_2[x]$. By Theorem~\ref{thm:sg-exact},
\[
  \SG_1(f_1) = 1 \oplus 2 \oplus 3 = 0, \qquad
  \SG_1(x \cdot f_1) = \SG_1(x^2 + x^3 + x^4) = 2 \oplus 3 \oplus 4 = 5.
\]
Then $\iota(\SG_1(x \cdot f_1)) = \Phi(x \cdot f_1) = \Phi(x) \boxtimes \Phi(f_1)
= \Phi(x) \boxtimes 0_R = 0_R = \iota(0)$, forcing $5 = \SG_1(x \cdot f_1) = 0$, a
contradiction.
\end{proof}

\begin{remark}[Characteristic-$2$ obstruction into $\Z$]
Into the standard ring of integers $(\Z,+,\cdot)$ (without modulo-$2$ or XOR addition), no
unital additive homomorphism $\Phi : \Z_2[x] \to \Z$ exists at all, since $1 + 1 = 0$ in
$\Z_2[x]$ forces $0 = \Phi(0) = \Phi(1+1) = \Phi(1) + \Phi(1) = 2$ in $\Z$.
\end{remark}

\subsection{Zero-Conductor and Quotient-Factorization Obstruction}

Finally, in algebraic coding theory and periodic coin-turning games, one sometimes asks
whether an SG function is ``generated by a homomorphism $\Phi : \Z_2[x] \to R$'' in the
weaker sense of \emph{factoring} as $\SG = \psi \circ \Phi$ through a quotient ring
$R = \Z_2[x]/\mathfrak{a}$ followed by an arbitrary set-theoretic readout $\psi : R \to \N$.
We prove that the largest ideal of $\Z_2[x]$ contained in $\ker(\SG_0)$ or $\ker(\SG_1)$
is $\{0\}$, ruling out every nontrivial quotient ring.

\begin{theorem}[Zero Conductor and Factorization Injectivity]
\label{thm:zero-conductor}
\begin{enumerate}
  \setlength{\itemsep}{2pt}
  \item For every nonzero polynomial $m(x) \in \Z_2[x]$, there exist integers $k_0 \ge 0$
    and $k_1 \ge 1$ such that $\SG_0(x^{k_0} m) \neq 0$ and $\SG_1(x^{k_1} m) \neq 0$.
  \item Let $(R,\boxplus,\boxtimes,0_R)$ be any rng, $\Phi : \Z_2[x] \to R$ any rng
    homomorphism, and $\psi : R \to \N$ any set map (with no algebraic assumptions on
    $\psi$). If either $\SG_0 = \psi \circ \Phi$ or $\SG_1 = \psi \circ \Phi$, then
    $\ker(\Phi) = \{0\}$ and $\Phi$ is injective.
\end{enumerate}
\end{theorem}

\begin{proof}
(1) Since $m \neq 0$, its support $\supp(m)$ is nonempty; let $d = \max\supp(m) = \deg(m) \ge 0$,
and set $r = d + 2$. Because $d + 2 < 2^{d+2} = 2^r$, both $k_0 = 2^r - 1 - d$ and
$k_1 = 2^r - d$ satisfy $k_0 \ge 1$ and $k_1 \ge 2$. Notice that
$\supp(x^k m) = \{j + k : j \in \supp(m)\}$.
\begin{itemize}
  \setlength{\itemsep}{2pt}
  \item For $\SG_0(x^{k_0} m) = \bigoplus_{j\in\supp(m)} (j + k_0 + 1)$: at $j = d$,
    $d + k_0 + 1 = 2^r$, whose $r$-th binary bit is $1$. For every $j \in \supp(m)$ with
    $j \neq d$, we have $j < d$ and hence $1 \le j + k_0 + 1 < d + k_0 + 1 = 2^r$, so its
    $r$-th binary bit is $0$. Thus $j = d$ is the \emph{unique} summand with bit $r$ equal
    to $1$, so bit $r$ of $\SG_0(x^{k_0} m)$ is $1$; in particular
    $\SG_0(x^{k_0} m) \ge 2^r > 0$.
  \item For $\SG_1(x^{k_1} m) = \bigoplus_{j\in\supp(m)} (j + k_1)$: at $j = d$,
    $d + k_1 = 2^r$ has bit $r$ equal to $1$, whereas every $j \in \supp(m)\setminus\{d\}$
    satisfies $1 \le j + k_1 < 2^r$ and has bit $r$ equal to $0$. Thus bit $r$ of
    $\SG_1(x^{k_1} m)$ is $1$, so $\SG_1(x^{k_1} m) \ge 2^r > 0$.
\end{itemize}

\smallskip
\noindent(2) Suppose $\SG_t = \psi \circ \Phi$ ($t \in \{0,1\}$), and let $m \in \ker(\Phi)$,
so $\Phi(m) = 0_R$. For every $k \ge 0$,
$\Phi(x^k m) = \Phi(x^k) \boxtimes \Phi(m) = \Phi(x^k) \boxtimes 0_R = 0_R = \Phi(0)$,
whence
\[
  \SG_t(x^k m) = \psi\bigl(\Phi(x^k m)\bigr) = \psi\bigl(\Phi(0)\bigr) = \SG_t(0) = 0.
\]
By part~(1), this forces $m = 0$; hence $\ker(\Phi) = \{0\}$. Finally, if
$\Phi(P) = \Phi(Q)$, then in $\Z_2[x]$ we have $Q + Q = 0$, so
$\Phi(P + Q) = \Phi(P) \boxplus \Phi(Q) = \Phi(Q) \boxplus \Phi(Q) = \Phi(Q + Q)
= \Phi(0) = 0_R$, which implies $P + Q \in \ker(\Phi) = \{0\}$ and therefore
$P = P + (Q + Q) = (P + Q) + Q = 0 + Q = Q$.
\end{proof}

\begin{corollary}[Impossibility over Algebraic $\F_2$-Algebras and Quotient Rings]
\label{cor:algebraic}
If $\Phi : \Z_2[x] \to R$ is any rng homomorphism with $\ker(\Phi) \neq \{0\}$---in
particular, any homomorphism onto a quotient ring $\Z_2[x]/(m)$ ($m \neq 0$), into a finite
ring, or into an algebraic $\F_2$-algebra such as Conway's nimber field
$(\N,\oplus,\otimes) \cong \overline{\F_2}$---there exists \emph{no} function
$\psi : R \to \N$ satisfying $\SG_0 = \psi \circ \Phi$ or $\SG_1 = \psi \circ \Phi$.
\end{corollary}

\section{Comparison of Coin-Turning Games and Formal Verification}

Table~\ref{tab:games} places Turning Turtles alongside related coin-turning games from
\emph{Winning Ways}~\cite{WW}, clarifying precisely when polynomial ring homomorphisms do
and do not arise.

\begin{table}[ht]
\centering
\caption{Algebraic status of Sprague--Grundy maps on $\Z_2[x]$ across coin-turning games.}
\label{tab:games}
\medskip
{\small
\begin{tabular}{llll}
\hline
Game & Single-coin SG & $\F_2$-linear? & Ring hom.\ to $(\N,\oplus,\otimes)$? \\
\hline
Single-Coin Flip & $1$ ($n \ge 0$) & Yes & \textbf{Yes} ($\ev_1$, $x \mapsto 1$) \\
Bounded Turtles (jump $\le 2$) & $2^{\otimes n} \in \{1,2,3\}$ & Yes & \textbf{Yes} ($\ev_2$, $x \mapsto 2$) \\
Turning Turtles ($0$-indexed) & $n + 1$ ($n \ge 0$) & Yes & \textbf{No} (fails at $x^3$: $4 \neq 1$) \\
Turning Turtles ($1$-indexed) & $n$ ($n \ge 1$) & Yes & \textbf{No} (fails at $x^2$: $2 \neq 1$) \\
Mock Turtles ($0$-indexed) & $2n$ or $2n{+}1$ (odious) & Yes & \textbf{No} (unbounded; $\ker$ not ideal) \\
\hline
\end{tabular}
}
\end{table}

\paragraph{Machine-Checked Formalization in Core Lean 4.19.0.}
The entire development in Sections~2--4 is formalized from scratch in
\texttt{lean/Main.lean} under Lean~4 toolchain \texttt{v4.19.0} with zero external
dependencies and strict warning-as-error checking enabled. Table~\ref{tab:lean} maps each
mathematical result to its formal Lean~4 declaration. All proofs rely solely on the standard
foundational axioms \texttt{propext}, \texttt{Classical.choice}, and \texttt{Quot.sound},
with zero \texttt{sorry}, \texttt{admit}, \texttt{native\_decide}, or custom axioms. An
independent Python verification script (\texttt{verify.py}) exhaustively cross-checks the
game-tree $\mex$ values, Conway nim-multiplication tables on $[0,16)$, the degree-$\le 2$
coincidence, the non-ideal kernel witnesses, and the zero-conductor bit-shift construction.

\begin{table}[ht]
\centering
\caption{Correspondence between mathematical theorems and Lean~4.19.0 declarations.}
\label{tab:lean}
\medskip
{\small
\begin{tabular}{>{\raggedright\arraybackslash}p{0.40\textwidth}>{\raggedright\arraybackslash}p{0.54\textwidth}}
\hline
Mathematical Result & Lean 4 Declaration in \texttt{Conjecture00000001288} \\
\hline
$\Z_2[x]$ Cauchy product \& ring laws &
  \texttt{PolyF2.coeff\_mul}, \texttt{PolyF2.mul\_assoc} \\
Well-founded game termination &
  \texttt{Move0.bits\_lt}, \texttt{Move1.bits\_lt} \\
Exact $\mex$ SG formulas (Thm.~\ref{thm:sg-exact}) &
  \texttt{sg0\_eq\_xorSum0}, \texttt{sg1\_eq\_xorSum1} \\
Surjective $\F_2$-linearity (Prop.~\ref{prop:linear}) &
  \texttt{sg0\_is\_F2\_linear\_surjective},\newline
  \texttt{sg1\_is\_F2\_linear\_surjective} \\
Unital nimber hom.\ $\ev_2$ &
  \texttt{ev2\_is\_nimber\_ring\_hom} \\
Deg-$\le 2$ match \& $x^3$ failure (Prop.~\ref{prop:coincidence}) &
  \texttt{sg0\_coincides\_with\_ev2\_deg\_le\_two\_}\allowbreak\texttt{and\_breaks\_at\_three} \\
Nimber ring refutation (Thm.~\ref{thm:nimber-no}) &
  \texttt{not\_nimber\_rng\_hom\_sg0},\newline
  \texttt{not\_nimber\_rng\_hom\_sg1} \\
Universal rng obstruction (Thm.~\ref{thm:universal-rng}) &
  \texttt{no\_rng\_hom\_generates\_sg0},\newline
  \texttt{no\_pos\_ideal\_rng\_hom\_generates\_sg1} \\
Zero conductor \& injectivity (Thm.~\ref{thm:zero-conductor}) &
  \texttt{sg0\_zero\_conductor},\newline
  \texttt{sg0\_factorization\_injective},\newline
  \texttt{sg1\_factorization\_injective} \\
Master resolution theorem &
  \texttt{conjecture\_00000001288\_resolution} \\
\hline
\end{tabular}
}
\end{table}

\begin{thebibliography}{9}

\bibitem{WW}
E.~R. Berlekamp, J.~H. Conway, and R.~K. Guy,
\emph{Winning Ways for Your Mathematical Plays}, Vol.~3, 2nd~ed.,
A~K Peters, Natick, MA, 2003, Chapter~14 (``Turn and Turn About''), pp.~461--494.

\bibitem{ONAG}
J.~H. Conway,
\emph{On Numbers and Games}, 2nd~ed.,
A~K Peters, Natick, MA, 2001, Chapter~6 (``The Curious Field $\mathbf{On}_2$''), pp.~50--67.

\bibitem{Lenstra}
H.~W. Lenstra, Jr.,
\emph{Nim multiplication},
S\'eminaire de Th\'eorie des Nombres de Bordeaux (1977--1978), Exp.~No.~11, 1--23.

\end{thebibliography}

\end{document}
